% !TeX program = XeLaTeX
% !TeX root = ../pujavidhanam.tex
\chapt{श्री-उमा-महेश्वर-पूजा}

\input{purvanga/vighneshwara-puja}

\sect{प्रधान-पूजा — श्री-उमा-महेश्वर-पूजा}

\shuklambaradharam

\pranayama

% TODO: set the occasion-specific calendrical values below
% \renewcommand{\ayane}{}
% \renewcommand{\rtu}{}
% \renewcommand{\masa}{}
% \renewcommand{\paksha}{}
% \renewcommand{\tithau}{}
% \renewcommand{\kaale}{}
\renewcommand{\prityartham}{श्री-उमामहेश्वर-प्रसादसिद्ध्यर्थं}
\renewcommand{\pujaam}{श्री-उमामहेश्वरपूजां}

\sankalpa{}

\vighneshvaraYathasthanam

\input{purvanga/aasana-puja}

\input{purvanga/ghanta-puja}

\input{purvanga/kalasha-puja}

\input{purvanga/aatma-puja}

\input{purvanga/pitha-puja}

\input{purvanga/guru-dhyanam}

\sect{षोडशोपचार-पूजा}
\renewcommand{\devAya}{ॐ उमामहेश्वराभ्यां नमः}
\begin{center}

\twolineshloka*
{मुक्तामालापरीताङ्गं दुकूलपरिवेष्टितम्}
{पञ्चाननमुमाकान्तमनन्तेन्दुरविप्रभम्}

\twolineshloka*
{चन्द्रार्धशेखरं नित्यं जटामकुटमण्डितम्}
{त्रिपुण्ड्ररेखाविलसद्भालभागमनामयम्}

\twolineshloka*
{भस्मोद्धूलितसर्वाङ्गं रुद्राक्षाभरणान्वितम्}
{मन्दस्मितमुखाम्भोजमाधारं जगतां प्रभुम्}

\twolineshloka*
{वेदैरनन्तैरनिशं स्तूयमानमनेकधा}
{वेदात्मकं वेदवेद्यं विष्णुब्रह्मादिवन्दितम्}

\twolineshloka*
{विष्णुनेत्रार्चितं भक्तकल्पद्रुममुमापतिम्}
{अप्रतिद्वन्द्वमद्वन्द्वं सर्ववृन्दारकार्चितम्}

\twolineshloka*
{सर्वोत्तममनाद्यन्तं सर्वदेवनिषेवितम्}
{ऋतं सत्यं परं ब्रह्म पुरुषं कृष्णपिङ्गलम्}

\twolineshloka*
{ऊर्ध्वरेतं विरूपाक्षं विश्वरूपं चिदात्मकम्}
{निष्कलं निर्गुणं शान्तं निरवद्यं निरञ्जनम्}

\twolineshloka*
{अप्रमेयं जगज्जन्मस्थितिसंहारकारणम्}
{विश्वबाहुं विश्वपादं विश्वाक्षं विश्वशम्भुवम्}

\twolineshloka*
{विश्वं नारायणाराध्यमक्षरं परमं पदम्}
{विश्वतः परमं नित्यं विश्वस्येशमनामयम्}

\medskip

\twolineshloka*
{लक्ष्मीसेवितपादाब्जां शचीसेवितपादुकाम्}
{सरस्वत्यादिभिर्नित्यं स्तूयमानपदाम्बुजाम्}

\twolineshloka*
{अधरोष्ठाधरीभूतपक्वबिम्बफलामिमाम्}
{मुखकान्तिकलोद्धूतपूर्णचन्द्रां शिवात्मिकाम्}

\twolineshloka*
{तिरस्कृतालिमालां तामलकावलिभिः सदा}
{पीनवक्षोजनिर्धूतचक्रवाकवराङ्गनाम्}

\twolineshloka*
{नित्यां तिरस्कृताम्भोजनयनां विश्वमातरम्}
{सीमन्तधिक्कृताशेषकामभल्लामहर्निशम्}

\twolineshloka*
{भ्रुकुटीनिर्जितानङ्गचापां भक्तजनेष्टदाम्}
{बाहुनालप्रभोद्धतहेमनालां विलासिनीम्}

\twolineshloka*
{रोमराजितिरोभूतसम्भ्रमभ्रमरावलिम्}
{नाभिरन्ध्रतिरोभूतघनावर्तोमिवर्तुलाम्}

\twolineshloka*
{उत्तमोरुतिरोभूतरम्भास्तम्भां शुभावहाम्}
{पादयुग्मप्रभापूरनिर्जितारुणपङ्कजाम्}

\twolineshloka*
{ब्रह्मेन्द्रोपेन्द्रजननीं महेशार्धाङ्गभागिनीम्}
{महेशाश्लिष्टवामाङ्गां वरदाभयदां सदा}

\twolineshloka*
{प्रसन्नवदनाम्भोजां स्मितपूर्वाभिभाषिणीम्}
{सर्वशृङ्गारवेषाढ्यां नानाभरणभूषिताम्}

देवैर्वेदैश्च वेदान्तैः स्तूयमानात्मवैभवाम्।
\textbf{ध्यायामि। \devAya{}।}
\medskip

\twolineshloka*
{महादेव कृपासिन्धो विश्ववन्द्यपदाम्बुज}
{आवाहयामि देवेश प्रसीद परमेश्वर}

\twolineshloka*
{लक्ष्म्यादिदेववनितापरिषेवितपादुके}
{आवाहयामि देवेशि प्रसीद परमेश्वरि}
\textbf{उमामहेश्वरौ इहागच्छतं इह तिष्ठतम्।}
\medskip

\twolineshloka*
{गृहाण सोम विश्वात्मन्नासनं रत्ननिर्मितम्}
{अनन्तासन विश्वेश करुणासागर प्रभो}

\twolineshloka*
{उमे सोमसमाश्लिष्टे सोमार्धकृतशेखरे}
{नानारत्नसमाकीर्णमासनं प्रतिगृह्यताम्}
\textbf{इदमासनम्। \devAya{}।}
\medskip

\twolineshloka*
{अर्घ्यं गृहाण देवेश सपुष्पं गन्धसंयुतम्}
{शिवानन्तगुणग्राम सर्वाभीष्टप्रदायक}

\twolineshloka*
{गृहाणार्घ्यं शिवे नित्यं सर्वावयवशोभिनि}
{शिवप्रिये शिवाकारे नित्यं भक्तवरप्रदे}
\textbf{हस्तयोः इदम् अर्घ्यम्। \devAya{}।}
\medskip

\twolineshloka*
{पाद्यं गृहाण गौरीश हेमपात्रस्थितं नवम्}
{गृहाण देवि पाद्यं त्वं वेदवेदान्तसंस्तुते}
\textbf{पादयोः इदं पाद्यम्। \devAya{}।}
\medskip

\twolineshloka*
{गृहाणाचमनं शम्भो कल्पितं निर्मल प्रभो}
{निर्मले निर्मलाकारे गृहाणाचमनीयकम्}
\textbf{मुखे इदम् आचमनीयम्। \devAya{}।}
\medskip

\twolineshloka*
{मधुपर्कं गृहाणेश सच्चिदानन्दविग्रह}
{मधुपर्कप्रदानेन प्रीतो भव महेश्वर}

\twolineshloka*
{मधुपर्कमिदं देवि स्वीकुरु प्रियशङ्करे}
{मधुपर्कप्रदानेन प्रीता भव सुशोभने}
\textbf{अयं मधुपर्कः। \devAya{}।}
\medskip

\twolineshloka*
{शिव शम्भोऽमृतस्नानं स्वीकुरुष्व कृपानिधे}
{सर्वतीर्थमयस्वामिन् सर्वपावनपावन}

\twolineshloka*
{शिवे पञ्चामृतस्नानं स्वीकृरुष्व शिवप्रिये}
{सृष्टतीर्थोत्तमे शुद्धे तीर्थराजनिषेविते}
\textbf{इदं पञ्चामृतस्नानम्, \devAya{}।}
\medskip

\twolineshloka*
{शम्भो शुद्धोदकस्नानं स्वीकुरुष्व सुरोत्तम}
{प्रसीद कृपया भक्तं पाहि मां पार्वतीपते}

\twolineshloka*
{शिवे शुद्धोदकस्नानं गृह्यतां मीनलोचने}
{प्रसीद देवि दीनं त्वं पाहि मां शरणागतम्}
\textbf{इदं शुद्धोदकस्नानम्। \devAya{}।}
\medskip

\twolineshloka*
{सोत्तरींयं सङ्गृहाण दुकूलमिदमुत्तमम्}
{पाहि पाहि कृपासिन्धो करुणावरुणालय}

\twolineshloka*
{सोत्तरीयं गृहाणेदं दुकूलं शङ्करप्रिये}
{प्रसीद पाहि मां दीनमनन्यशरणं शिवे}
\textbf{इमे वस्त्रोपवस्त्रे। \devAya{}।}
\medskip

\twolineshloka*
{उपवीतं गृहाणेश शम्भो चन्द्रार्धशेखर}
{गृह्यतामुपवीतं यन्नवतन्तुविनिर्मितम्}
\textbf{इमे यज्ञोपवीते। \devAya{}।}
\medskip

\twolineshloka*
{प्रसीद पार्वतीनाथ शरणागतवत्सल}
{गृहाण चन्दनं देवि कस्तूरीकुङ्कुमाञ्चितम्}

\twolineshloka*
{देवि देवाधिदेवेशि यथाशक्त्युपकल्पितम्}
{गृहाण चन्दनं देव नानागन्धोपकल्पितम्}
\textbf{अयं गन्धः, \devAya{}।}
\medskip

\twolineshloka*
{विश्वाभरण विश्वेश रत्नाभरणभूषित}
{गृहाणाभरणानीश निरीश निगमाश्रय}

\twolineshloka*
{विश्वं विश्वात्मिके पाहि विश्वनाथप्रिये सदा}
{गृहाणाभरणान्यम्ब सर्वाभरणभूषिते}
\textbf{इदं आभूषणं \devAya{}।}
\medskip

\twolineshloka*
{सर्वप्रिये जगद्वन्द्ये जगदानन्ददे शिवे}
{गृहाण बिल्वपत्राणि पुष्पाणि च महेश्वर}
\textbf{इमानि पुष्पाणि। ॐ उमामहेश्वराभ्यां नमः।}
\medskip

\twolineshloka*
{सुगन्धीनि नवानीश सबिल्वकुसुमप्रिय}
{गृहाण बिल्वपत्राणि नूतनानि शिवप्रिये}
\textbf{एतानि बिल्वपत्राणि \devAya{}।}
\medskip

\end{center}
\begingroup
\centering
\setlength{\columnseprule}{1pt}
\let\chapt\sect
\input{../namavali-manjari/100/Ardhanarishvara_108.tex}
\endgroup


\begin{center}

\twolineshloka*
{नागपुन्नागमन्दारमालिकासमलङ्कृते}
{दशाङ्गं गुग्गुलुं धूपं सगोघृतमनुत्तमम्}

\twolineshloka*
{गृहाण पार्वतीनाथ घ्राणतर्पणमादरात्}
{दशाङ्गं गुग्गुलुं धूपं सगोघृतमनुत्तमम्}
\textbf{अयं धूपः। \devAya{}।}
\medskip

\twolineshloka*
{गृहाण भक्तवरदे लक्ष्मीवाण्यादिसेविते}
{साज्यं त्रिवर्तिविस्फूर्जद्दीपं शिव शिवापते}

\twolineshloka*
{गृहाणानन्तसूर्याग्निचन्द्रप्रभ नमोऽस्तु ते}
{साज्यं त्रिवर्तिविस्पूर्जद्दीपमीशानवल्लभे}
\textbf{एष दीपः \devAya{}।}
\medskip

\twolineshloka*
{गृहाण चन्द्रसूर्याग्निमण्डलाधिकसुप्रभे}
{शम्भो गोघृतसंयुक्तं परमान्नं मनोहरम्}

\twolineshloka*
{सशर्करं गृहाणेश नित्यतृप्त महेश्वर}
{शिवे गोवृतसंयुक्तं परमान्नं मनोहरम्}
\textbf{एतानि नानाविध-नैवेद्यानि \devAya{}।}
\medskip

\twolineshloka*
{सशर्करं गृहाणाम्ब परमान्नप्रदायिनि}
{शम्भो गृहाण गन्धाढ्यमिदमाचमनीयकम्}

\twolineshloka*
{शुद्धे शुद्धिप्रदे देवि शिववामाङ्गरूपिणि}
{शिवे गृहाण गन्धाढ्यमिदमाचमनीयकम्}
\textbf{इदं आचमनीयं \devAya{}।}
\medskip

\twolineshloka*
{नीराजनं गृहाणेश बहुदीपविरजितम्}
{स्वप्रकाश प्रकाशात्मन्प्रकाशितदिगन्तर}

\twolineshloka*
{नीराजनं गृहाणाम्ब सूर्यचन्द्राग्निलोचने}
{प्रभापूरित सर्वाङ्गे मङ्गले मङ्गलास्पदे}
\textbf{इदं नीराजनं \devAya{}।}
\medskip

\twolineshloka*
{शम्भो गृहाण ताम्बूलमेलाकर्पूरसंयुतम्}
{प्रसीद भगवन् शम्भो सर्वज्ञामितविक्रम}

\twolineshloka*
{शिवे गृहाण ताम्बूलमेलाकर्पूरसंयुतम्}
{प्रसीद सस्मिते देवि शिवालिङ्गितविग्रहे}
\textbf{इदं ताम्बूलम्। \devAya{}।}
\medskip

\twolineshloka*
{गृहाण परमेशान सरत्ने छत्रचामरे}
{दर्पणं व्यजनं चेश सर्वदुःखविनाशक}

\twolineshloka*
{गृहाणेमे सुराराध्ये सरत्ने छत्रचामरे}
{दर्पणं व्यजनं चाद्ये विद्याधारे नमोऽस्तु ते}
\textbf{इदं छत्रं, इदं चामरं, इदं व्यजनम्, इदं दर्पणः। \devAya{}।}
\medskip

\twolineshloka*
{प्रदक्षिणनमस्कारान्गृहाण परमेश्वर}
{नर्तनं च महादेव नाट्यप्रिय नमोऽस्तु ते}

\twolineshloka*
{प्रदक्षिणनमस्कारान्गृहाण परमेश्वरि}
{नर्तनं च महादेवि शिवनाट्यम्प्रिये शिवे}
\medskip

अनया पूजया भगवान् उमामहेश्वरः प्रीयताम्॥

\end{center}

\sect{श्री-उमामहेश्वराष्टकम्}
\begin{center}

\twolineshloka*
{पितामहशिरच्छेदप्रवीणकरपल्लव}
{नमस्तुभ्यं नमस्तुभ्यं नमस्तुभ्यं महेश्वर}

\twolineshloka*
{निशुम्भशुम्भप्रमुखदैत्यशिक्षणदक्षिणे}
{नमस्तुभ्यं नमस्तुभ्यं नमस्तुभ्यं महेश्वरि}

\twolineshloka*
{शैलराजस्य जामातः शशिरेखावतंसक}
{नमस्तुभ्यं नमस्तुभ्यं नमस्तुभ्यं महेश्वर}

\twolineshloka*
{शैलराजात्मजे मातः शातकुम्भनिभप्रभे}
{नमस्तुभ्यं नमस्तुभ्यं नमस्तुभ्यं महेश्वरि}

\twolineshloka*
{भूतनाथपुराराते भुजङ्गकृतभूषण}
{नमस्तुभ्यं नमस्तुभ्यं नमस्तुभ्यं महेश्वर}

\twolineshloka*
{पादप्रणतभक्तानां पारिजातगुणाधिके}
{नमस्तुभ्यं नमस्तुभ्यं नमस्तुभ्यं महेश्वरि}

\twolineshloka*
{हालास्येश दयामूर्त्ते हालाहललसद्गल}
{नमस्तुभ्यं नमस्तुभ्यं नमस्तुभ्यं महेश्वर}

\twolineshloka*
{नितम्बिनि महेशस्य कदम्बवननायिके}
{नमस्तुभ्यं नमस्तुभ्यं नमस्तुभ्यं महेश्वरि}

इति श्रीहालास्यमाहात्म्ये सङ्घिलकृतम् उमामहेश्वराष्टकं सम्पूर्णम्।

\end{center}

\closesection
